\documentclass[a4paper,10pt,ngerman]{INSOproject}

\dokumenttyp{Projektbericht}
\assistent{Vorname Nachname}
\betreuer{Thomas Grechenig}
\title{Weiterentwicklung einer Serious-Game-Plattform}
\shorttitle{Projektbericht Serious-Game-Plattform}
\subtitle{Projektstand, Beitrag im Vier-Monats-Zeitraum und Android-Launcher-Kontext}
\author{Vorname Nachname}
\matrikelnr{Matrikelnummer}
\kennzahl{Kennzahl}
\studium{Studienrichtung}
\date{\today}

\bibliography{db}

\begin{document}

\maketitle
\tableofcontents
\newpage

\section{Problemstellung}
Digitale Anwendungen gewinnen im Gesundheitskontext zunehmend an Bedeutung. Neben klassischen Web- und Mobile-Anwendungen entstehen immer mehr interaktive Systeme, die Training, Motivation und Verlaufskontrolle miteinander verbinden. In diesem Umfeld nehmen \textbf{Serious Games} eine besondere Rolle ein, weil sie spielerische Elemente mit funktionalen Zielen wie Rehabilitation, Aktivierung, Therapiebegleitung oder strukturierter Datenerfassung verknüpfen.

Das vorliegende Projekt befasst sich mit der Weiterentwicklung einer Serious-Game-Plattform, die mehrere technische Bausteine in einem gemeinsamen System zusammenführt. Dazu gehören ein Backend zur Daten- und Benutzerverwaltung, eine Weboberfläche für administrative und therapeutische Prozesse, ein Android-Launcher zur Bereitstellung mobiler Spiele sowie eine Test-App zur Überprüfung der Session- und Score-Übertragung.

Die Problemstellung ergibt sich daraus, dass Serious Games in der Praxis nicht als isolierte Einzelanwendungen betrachtet werden können. Damit sie im Gesundheitskontext sinnvoll einsetzbar sind, muss eine Plattform mehrere Anforderungen gleichzeitig erfüllen:

\begin{itemize}
  \item Spiele müssen zentral verwaltet und mit konsistenten Metadaten beschrieben werden.
  \item Relevante Spiele müssen bestimmten Nutzungskontexten oder Patientinnen und Patienten zugeordnet werden können.
  \item Mobile Spiele müssen auf dem Zielgerät zuverlässig verfügbar sein und sich dort starten lassen.
  \item Spielsessions, Scores und weitere Ergebnisdaten müssen an die Plattform zurückgeführt und sichtbar gemacht werden.
\end{itemize}

Die zentrale Herausforderung des Projekts liegt somit nicht in der Entwicklung eines einzelnen Spiels, sondern in der Schaffung einer belastbaren Plattformlogik für Verwaltung, Bereitstellung, Nutzung und Auswertung.

\subsection{Was ist das Projekt?}
Die Serious-Game-Plattform ist als verteiltes Softwaresystem angelegt. Sie verbindet mehrere Komponenten, die gemeinsam den Lebenszyklus eines Serious Games abbilden:

\begin{itemize}
  \item \textbf{Backend:} Persistenz, Rollenlogik, Upload-Prozesse, Zuweisungen und Session-Daten.
  \item \textbf{Webanwendung:} Verwaltungsoberfläche für Admin- und Prozessschritte rund um Serious Games.
  \item \textbf{Android-Launcher:} Mobile Anwendung zur Anzeige und zum Start zugewiesener Android-Spiele.
  \item \textbf{Test-App:} Hilfsanwendung zur technischen Überprüfung der Session- und Score-Übertragung.
\end{itemize}

Das Projekt sagt fachlich aus, dass Serious Games im Gesundheitskontext erst dann praktisch nutzbar werden, wenn diese Systemteile zuverlässig zusammenspielen. Eine Plattform muss deshalb nicht nur Inhalte speichern, sondern den vollständigen Pfad vom Upload eines Spiels bis zur Anzeige seiner Ergebnisse abbilden.

\subsection{Aktueller Stand des Projekts}
Zum aktuellen Stand liegt eine Plattform vor, deren Kernbausteine funktional miteinander verbunden sind. Besonders relevant ist die inzwischen bessere Verzahnung zwischen Upload-Prozess, mobiler Bereitstellung und Sichtbarkeit von Spielergebnissen.

Der derzeitige Projektstand lässt sich wie folgt zusammenfassen:

\begin{itemize}
  \item Es existiert eine zentrale Struktur zur Verwaltung von Serious Games und zugehörigen Metadaten.
  \item Mobile Spiele können in den Plattformkontext integriert werden, insbesondere über APK-basierte Prozesse.
  \item Der Android-Launcher bildet die patientennahe Ausführungsebene für mobile Spiele.
  \item Sessions und Scores können an die Plattform zurückgegeben und in der Oberfläche sichtbar gemacht werden.
  \item Für Integrations- und Funktionsprüfungen steht eine Test-App als technisches Hilfsmittel zur Verfügung.
\end{itemize}

Aus Projektsicht ist die Lösung damit nicht mehr nur ein loser Prototypenverbund, sondern besitzt bereits einen zusammenhängenden End-to-End-Charakter. Gleichzeitig bleibt sie ausbaufähig, etwa bei der weiteren Automatisierung, bei robusterer Metadatenanalyse und bei der Testabdeckung.

\section{Methodisches Vorgehen}
Die Arbeit am Projekt erfolgte schrittweise entlang der bestehenden Architektur. Ziel war es, zunächst die vorhandenen Zusammenhänge sauber zu verstehen und darauf aufbauend die kritischen Integrationsstellen weiterzuentwickeln.

\subsection{Monat 1: Einarbeitung und Systemverständnis}
Im ersten Monat lag der Schwerpunkt auf dem Verständnis der bestehenden Plattform. Dafür wurden die vorhandenen Komponenten, ihre Verantwortlichkeiten und ihre Schnittstellen zueinander analysiert. Ziel war es, die Plattform nicht nur auf Dateiebene zu verstehen, sondern den tatsächlichen Ablauf von Spielverwaltung, Zuweisung, Bereitstellung und Nutzung nachzuvollziehen.

Schwerpunkte dieser Phase:
\begin{itemize}
  \item Analyse der Architektur aus Backend, Webanwendung, Android-Komponenten und Testpfaden.
  \item Nachvollzug des Datenflusses zwischen Upload, Assignment, Session-Übertragung und Anzeige.
  \item Identifikation manueller und potenziell fehleranfälliger Schritte im mobilen Bereitstellungsprozess.
  \item Aufbau einer belastbaren Grundlage für spätere Änderungen über mehrere Systemteile hinweg.
\end{itemize}

\subsection{Monat 2: Backend- und Integrationslogik}
Im zweiten Monat stand die technische Weiterentwicklung der Integrationslogik im Vordergrund. Dabei ging es insbesondere um Abläufe, die mobile Spiele in den Plattformkontext einbinden und deren Daten später wieder verarbeitbar machen.

Wesentliche Tätigkeiten:
\begin{itemize}
  \item Überprüfung und Stabilisierung von API-nahen Abläufen im Serious-Game-Bereich.
  \item Vereinheitlichung von Validierungs- und Fehlerpfaden für relevante Artefakte.
  \item Verbesserung der Nachvollziehbarkeit im Umgang mit Spiel-Metadaten.
  \item Vorbereitung robusterer Integrationspfade für APK-basierte Spiele.
\end{itemize}

In dieser Phase war besonders wichtig, technische Änderungen so anzulegen, dass sie nicht nur lokal funktionieren, sondern in den bestehenden Plattformfluss passen.

\subsection{Monat 3: Mobile Bereitstellung und Android-Launcher-Kontext}
Im dritten Monat verlagerte sich der Fokus auf die mobile Seite des Projekts. Ziel war es, den Weg von einem eingebundenen mobilen Spiel bis zur sichtbaren Verfügbarkeit im Android-Launcher konsistenter und praxisnäher zu gestalten.

Dabei standen folgende Punkte im Mittelpunkt:
\begin{itemize}
  \item Verbesserung des Upload- und Zuordnungsflusses für mobile Spielpakete.
  \item Berücksichtigung von APK-spezifischen Eigenschaften innerhalb des Plattformmodells.
  \item Reduktion manueller Zwischenschritte bei der Bereitstellung im Launcher-Kontext.
  \item Bessere Grundlage für eine direkte Nutzung neuer mobiler Spiele auf dem Zielgerät.
\end{itemize}

Diese Phase war für den Projektcharakter besonders wichtig, weil hier sichtbar wurde, wie stark technische Infrastruktur und tatsächliche Nutzbarkeit zusammenhängen.

\subsection{Monat 4: Sichtbarkeit von Sessions, Scores und Testpfaden}
Im vierten Monat lag der Schwerpunkt auf der Frage, wie Ergebnisse aus Serious Games nicht nur technisch empfangen, sondern auch sinnvoll sichtbar gemacht werden können. Gleichzeitig wurden Testpfade genutzt, um die Integrationskette zwischen Spiel, Launcher, Test-App und Plattform besser nachvollziehen zu können.

Umgesetzt beziehungsweise vertieft wurden dabei:
\begin{itemize}
  \item Sichtbarmachung von Session-Daten in der Benutzeroberfläche.
  \item Einbindung beziehungsweise Nutzung von Score- und Highscore-Informationen.
  \item Verbesserung der Nachvollziehbarkeit von Session-Metadaten.
  \item Nutzung einer Test-App zur Validierung des technischen End-to-End-Verhaltens.
\end{itemize}

\section{Resultat}
Aus dem viermonatigen Projektzeitraum lassen sich mehrere konkrete Ergebnisse ableiten:

\begin{enumerate}
  \item \textbf{Klarerer Projektzusammenhang:} Die Plattform wurde als zusammenhängendes System aus Verwaltung, mobiler Bereitstellung und Auswertung weiter geschärft.
  \item \textbf{Verbesserte mobile Einbindung:} APK-bezogene Abläufe wurden stärker in den Plattformprozess integriert.
  \item \textbf{Bessere Sichtbarkeit von Ergebnissen:} Sessions und Scores sind für die Plattformnutzung relevanter und nachvollziehbarer geworden.
  \item \textbf{Unterstützung durch Testpfade:} Mit der Test-App kann das Zusammenspiel der Komponenten besser überprüft werden.
  \item \textbf{Stärkere Erweiterbarkeit:} Die Plattform wurde in eine Richtung weiterentwickelt, in der zusätzliche Spiele und mobile Szenarien einfacher anschließbar sind.
\end{enumerate}

\subsection{Einordnung der Ergebnisse}
Der aktuelle Projektstand zeigt, dass die Serious-Game-Plattform eine belastbare Grundlage für die Verwaltung, Bereitstellung und Auswertung digitaler Therapie- und Trainingsspiele bietet. Besonders relevant ist die stärkere Verbindung zwischen administrativen Prozessen, mobiler Nutzung über den Android-Launcher und der Rückführung von Session-Daten in die Plattform.

Innerhalb von vier Monaten lag mein Beitrag vor allem in der Analyse und Weiterentwicklung genau dieser Integrationspunkte. Dadurch wurde die Plattform in einem Bereich gestärkt, der für ihren praktischen Nutzen entscheidend ist: dem Übergang von einem verwalteten Spielobjekt zu einem tatsächlich nutzbaren und auswertbaren Serious Game.

\subsection{Herausforderungen}
Während des Projekts zeigten sich mehrere technische Herausforderungen:

\begin{itemize}
  \item \textbf{Systemgrenzen und Schnittstellen:} Schon kleine Änderungen an Datenmodellen oder Integrationslogik wirken sich auf Backend, Weboberfläche, Launcher und Testpfade aus.
  \item \textbf{Mobile Bereitstellung:} Die Einbindung von Android-Spielen ist anspruchsvoller als ein reiner Upload-Prozess, weil Metadaten, Zuweisungen und Sichtbarkeit im Launcher zusammenpassen müssen.
  \item \textbf{Nachvollziehbarkeit von Ergebnissen:} Session- und Score-Daten sind nur dann fachlich wertvoll, wenn sie konsistent übertragen, gespeichert und angezeigt werden.
\end{itemize}

\subsection{Ausblick}
Für weitere Ausbaustufen sind insbesondere folgende Themen sinnvoll:

\begin{itemize}
  \item robustere und stärker automatisierte Metadatenanalyse für mobile Spielpakete,
  \item erweiterte Auswertungsansichten für Sessions und Scores,
  \item zusätzliche Integrations- und End-to-End-Tests,
  \item klarere Deployment- und Betriebsprozesse für mobile Spiele im Plattformkontext.
\end{itemize}

\section{Platzhalter für Screenshots}
Die folgenden Abschnitte dienen als Platzhalter für spätere visuelle Ergänzungen im Projektbericht.

\subsection{Launcher Startansicht}
\fbox{\parbox[c][7cm][c]{\textwidth}{\centering Platzhalter: Android-Launcher Startansicht}}

\subsection{Launcher Spieleübersicht}
\fbox{\parbox[c][7cm][c]{\textwidth}{\centering Platzhalter: Launcher mit zugewiesenen Spielen}}

\subsection{Launcher Detail- oder Auswahlansicht}
\fbox{\parbox[c][7cm][c]{\textwidth}{\centering Platzhalter: Launcher-View mit Spielauswahl oder Detailansicht}}

\subsection{Webansicht für Spielverwaltung}
\fbox{\parbox[c][7cm][c]{\textwidth}{\centering Platzhalter: Webansicht für Serious-Game-Verwaltung oder Upload}}

\subsection{Session- und Score-Ansicht}
\fbox{\parbox[c][7cm][c]{\textwidth}{\centering Platzhalter: Anzeige von Sessions, Scores oder Highscores}}

\subsection{Test-App Startansicht}
\fbox{\parbox[c][7cm][c]{\textwidth}{\centering Platzhalter: Test-App Startansicht}}

\subsection{Test-App Session-Übertragung}
\fbox{\parbox[c][7cm][c]{\textwidth}{\centering Platzhalter: Test-App bei Session- oder Score-Submission}}

\appendix
\section{Anhang: Einordnung des Android-Launchers}
Der Android-Launcher ist innerhalb der Plattform die mobile Ausführungsebene für zugewiesene Spiele. Aus Projektsicht erfüllt er damit nicht nur die Funktion einer Startoberfläche, sondern bildet die Verbindung zwischen zentral verwaltetem Spielbestand und tatsächlicher Nutzung auf dem Endgerät.

Die Funktion des Launchers lässt sich für den Projektbericht wie folgt zusammenfassen:

\begin{itemize}
  \item Der Launcher stellt die auf dem Gerät verfügbaren beziehungsweise zugewiesenen Android-Spiele in einer nutzbaren Oberfläche dar.
  \item Er dient als patientennahe Einstiegskomponente und reduziert den Aufwand, einzelne Spiele manuell zu suchen oder getrennt zu starten.
  \item Neue oder aktualisierte mobile Spiele werden nicht isoliert betrachtet, sondern in einen Plattformprozess aus Verwaltung, Zuweisung und Nutzung eingebettet.
  \item Für den Projektkontext ist der Launcher deshalb ein zentrales Bindeglied zwischen Web-/Backend-Verwaltung und realem Einsatz des Serious Games auf Android.
\end{itemize}

\subsection{Bedeutung für den Projektbericht}
Für diesen Bericht ist der Android-Launcher vor allem deshalb relevant, weil an ihm sichtbar wird, ob die Plattform ihr Ziel tatsächlich erreicht. Erst wenn ein Spiel nach Upload, Zuordnung und technischer Integration im Launcher erscheint und dort gestartet werden kann, wird aus einer Verwaltungsfunktion ein praktisch nutzbarer Therapie- oder Trainingsbaustein.

\subsection{README-Auszug}
\fbox{\parbox[c][6cm][c]{\textwidth}{\centering Platzhalter: Hier kann bei Bedarf der konkrete README-Auszug zum Android-Launcher eingefügt werden.}}

\printbibliography

\end{document}
